\centerline{\bf \Huge Математическая индукция}
\centerline{\bf \Huge Знакомство}

\begin{flushright}
        \scriptsize{}
\end{flushright}

Туториал по шаганию по бесконечной лестнице:

1) \textbf{База.} Найдите начало лестницы. Первым делом нужно найти первую ступеньку и уверенно встать на неё.

2) \textbf{Переход.} После того как встали на первую ступеньку, нужно научиться перешагивать на следующую ступеньку. 

Повторяйте шаг номер 2 бесконечное количество раз. Поздравляю, вы поняли, что такое математическая индукция!

\problem{1}
Рассмотрим уголок. Он получается вырезанием из квадрата 2x2 одной клетки. Можно ли разрезать на такие уголки квадрат следующих размеров без одной клетки (вырезана может быть любая клетка квадрата, даже откуда-то из середины)

$a)4\times4 \quad b)8\times8  \quad c)128\times128 \quad d)2^n\times2^n ?$


\problem{2}
Из $3^n$ монет одна фальшивая, которая легче настоящей. Докажите, что за $n$
взвешиваний можно найти фальшивую монету

\problem{3}
В прямоугольнике $3 \times n$ расставлены фишки трех цветов по $n$ каждого цвета. Докажите, что переставляя фишки в строчках, можно сделать так, чтобы в каждом столбце все три фишки стали разными.

\problem{4}
На плоскости проведено несколько прямых. Докажите, что части, на которые эти
прямые делят плоскость, можно раскрасить в два цвета так, чтобы любые две
соседние части были окрашены в разные цвета.

\problem{5}
Докажите, что если $x + \dfrac{1}{x}$ целое, то $x^n +\dfrac{1}{x^n}$ тоже целое.

\problem{6}
Пятеро разбойников добыли мешок золотого песка. Они хотят поделить его так, чтобы каждый был уверен, что он получил не меньше одной пятой золота. Никаких способов измерения у них нет, однако каждый умеет оценивать на глаз величину кучи песка. Мнения разбойников о величине куч могут расходиться. Как им поделить добычу?
